\documentclass[11pt]{report}
\usepackage[margin=1in]{geometry}
\usepackage{amsmath,amssymb}
\usepackage{graphicx}
\graphicspath{{figs/}}
\usepackage{booktabs}
\usepackage{float}
\usepackage[hidelinks]{hyperref}

\title{A Physics-Informed Graph Neural Network Surrogate for a\\
Charring-Ablator Material Response Solver}
\author{Grey Goodwin}
\date{}

\begin{document}
\maketitle

\chapter{The Reference Solver}

The surrogate developed in this work is trained on data produced by a
one-dimensional finite-volume material-response solver for a charring ablator.
The solver belongs to a well-established family. The coupled treatment of a
decomposing solid and a chemically reacting boundary layer was set out in the
Aerotherm charring-material ablation analysis~\cite{kendall1968part1,moyer1968cma},
and that formulation still underlies current material-response codes such as
FIAT~\cite{chen1999fiat}, CHAR~\cite{amar2016char}, and
PATO~\cite{lachaud2014pato}. The particular solver used here shares its lineage
with the material-response work of Martin and
colleagues~\cite{martin2010nondarcian,weng2014pyrolysisgas,martin2015coupled}.

This chapter states the physics that solver implements. The purpose is twofold.
First, the equations below define what the surrogate is being asked to reproduce.
Second, several of them are later built directly into the surrogate rather than
learned, so they need to be stated precisely before that construction can be
described.

Every equation in this chapter is transcribed from the solver source rather than
from a textbook. Where the implemented form departs from the form usually written
in the literature, the implemented form is given and the departure is noted.

\section{Domain and Discretization}

The material is treated as a slab of thickness $L_0$ resolved by $N$ uniform
finite-volume cells in the standard control-volume
arrangement~\cite{patankar1980,versteeg2007fvm}. Cell centers are indexed
$i = 1 \dots N$. Two ghost cells,
indexed $0$ and $N+1$, sit outside the domain and carry the boundary conditions.
The aerothermally heated surface is the right end of the domain at $x = L_0$, and
the back face is the left end at $x = 0$.

The cell width, face area, and cell volume are
\begin{equation}
\Delta x = \frac{L_0}{N}, \qquad
dA_i = (2x_i)^{m}\,\pi^{m(3-m)/2}, \qquad
dV_i = \frac{2^{m}\pi^{m(3-m)/2}}{m+1}\left(x_{i+1}^{m+1} - x_i^{m+1}\right),
\end{equation}
where $m = 0$, $1$, or $2$ selects Cartesian, cylindrical, or spherical geometry.
All cases considered here use $m = 0$, so $dA_i = 1$ and $dV_i = \Delta x$.

Two geometric coefficients recur in the flux terms,
\begin{equation}
\mathrm{FC}^-_i = \frac{2}{\Delta x_i + \Delta x_{i-1}}, \qquad
\mathrm{FC}^+_i = \frac{2}{\Delta x_{i+1} + \Delta x_i},
\end{equation}
and a face-interpolation weight is used to reconstruct face values from cell
values,
\begin{equation}
\mathrm{g}^-_i = \frac{\Delta x_{i-1}}{\Delta x_{i-1} + \Delta x_i}, \qquad
\mathrm{g}^+_i = \frac{\Delta x_{i+1}}{\Delta x_i + \Delta x_{i+1}} .
\end{equation}

\section{Energy Equation}

The solver does not advance temperature directly. It advances the total energy
stored in each cell and then recovers temperature by inverting a nonlinear
enthalpy relation. This distinction matters, because the material's heat capacity
and its chemical state both change as it decomposes, and an enthalpy formulation
carries the energy released by decomposition without a separate heat-of-pyrolysis
term. The underlying in-depth energy balance for a decomposing ablator is the one
derived in the charring-material ablation analysis~\cite{moyer1968cma}.

The quantity conserved in each cell is
\begin{equation}
E_i = \underbrace{\rho_i\, h_s(T_i, \tau_i)\, dV_i}_{\text{solid sensible enthalpy}}
    + \underbrace{\varphi_i\, \rho_g(T_i, p_{g,i})\, u_g(T_i)\, dV_i}_{\text{gas internal energy}},
\label{eq:energy_state}
\end{equation}
where $\rho$ is the bulk solid density, $h_s$ the mass-specific sensible enthalpy
of the solid, $\tau$ the char-progress variable defined in
Section~\ref{sec:tau}, $\varphi$ the porosity, $\rho_g$ the pyrolysis-gas
density, and $u_g$ the gas specific internal energy. The gas term is present only
when gas-energy storage is enabled, which is the case for all runs used here.

Over one time step the cell energy changes by conduction across its two faces, by
advection of pyrolysis gas, and by the change in chemical state,
\begin{equation}
E_i^{n+1} = E_i^{n}
  + \Delta t \Big[ \underbrace{-F^-_i\, dA_i + F^+_i\, dA_{i+1}}_{\text{net conduction}}
  - \underbrace{\left(H_{i+1} - H_i\right)}_{\text{gas advection}} \Big]
  + E^{\mathrm{chem}}_i .
\label{eq:energy_update}
\end{equation}

The conduction fluxes are evaluated from temperature differences across faces
using the face conductivity defined below,
\begin{equation}
F^-_i = k_{f,i}\,\mathrm{FC}^-_i\,(T_i - T_{i-1}), \qquad
F^+_i = k_{f,i+1}\,\mathrm{FC}^+_i\,(T_{i+1} - T_i).
\end{equation}

The advected quantity is the enthalpy carried by the pyrolysis gas crossing each
face,
\begin{equation}
H_i = \dot m_i\, h_g(T_{f,i}), \qquad
T_{f,i} = T_{i-1} + (T_i - T_{i-1})\,\mathrm{g}^-_i ,
\end{equation}
where $\dot m_i$ is the gas mass flux at face $i$ from Section~\ref{sec:gas} and
$h_g$ is the gas enthalpy relative to the reference temperature. At an open
boundary the advected enthalpy is set to zero when the flux is directed inward,
so that ambient gas does not carry spurious enthalpy into the domain.

The chemical term is the change in the blended formation enthalpy as the solid
decomposes,
\begin{equation}
E^{\mathrm{chem}}_i = \Big[\rho_i^{n+1} h_{s,\mathrm{form}}(\tau_i^{n+1})
  - \rho_i^{n} h_{s,\mathrm{form}}(\tau_i^{n})\Big] dV_i .
\label{eq:echem}
\end{equation}
This is the only route by which pyrolysis energy enters the balance. There is no
explicit $\dot\rho\,\Delta h_{\mathrm{pyro}}$ source. The heat of pyrolysis is
implicit in the offset between the virgin and char enthalpy curves at the
reference temperature, and it is released as $\tau$ evolves.

In continuous form, Equations~\eqref{eq:energy_state} through \eqref{eq:echem}
discretize
\begin{equation}
\frac{\partial}{\partial t}\Big[\rho\,h_s(T,\tau) + \varphi\,\rho_g u_g\Big]
= \frac{1}{dV}\frac{\partial}{\partial x}\!\left(A\,k_{\mathrm{mix}}\frac{\partial T}{\partial x}\right)
- \frac{\partial}{\partial x}\!\left(\dot m''\, h_g\right)
- \frac{\partial}{\partial t}\Big[\rho\,h_{s,\mathrm{form}}(\tau)\Big].
\label{eq:energy_continuous}
\end{equation}

\subsection{Face Conductivity}

The conductivity at a face is not an average of the two neighboring cell
conductivities. The material state is interpolated to the face first, and the
conductivity function is then evaluated once at that interpolated state,
\begin{equation}
\rho_{f,i} = \tfrac12\left(\rho_{i-1} + \rho_i\right), \qquad
T_{f,i} = T_i + (T_{i-1} - T_i)\,\mathrm{g}^+_{i-1},
\end{equation}
\begin{equation}
k_{f,i} = \tau_{f,i}\, k_v(T_{f,i}) + \left(1 - \tau_{f,i}\right) k_c(T_{f,i}),
\qquad \tau_{f,i} = \tau(\rho_{f,i}).
\label{eq:kface}
\end{equation}
The density used in Equation~\eqref{eq:kface} is the density at the start of the
step, before the pyrolysis update. Using the updated density here biases the
conductivity in hot cells where decomposition is fast.

\subsection{Temperature Recovery}

Equation~\eqref{eq:energy_update} gives the new cell energy. Temperature is
recovered from it by a local Newton iteration, since $h_s$ depends nonlinearly on
$T$. A predictor is formed first,
\begin{equation}
T^{(0)}_i = T^n_i + \alpha_i\left(\text{net flux} - \text{advection}\right),
\qquad
\alpha_i = \frac{\Delta t / dV_i}{\rho^{n+1}_i c_{p,\mathrm{mix}}(T^n_i, \bar\tau_i)
  + \bar\varphi_i \rho_g c_{v,g}},
\end{equation}
and is then corrected by
\begin{equation}
T^{(m+1)}_i = T^{(m)}_i +
\frac{E^{n+1}_i - \left[\rho^{n+1}_i h_s(T^{(m)}_i, \tau^{n+1}_i)
  + \varphi^{n+1}_i \rho_g u_g\right] dV_i}
{\left[\rho^{n+1}_i c_{p,\mathrm{mix}}(T^{(m)}_i, \tau^{n+1}_i)
  + \varphi^{n+1}_i \rho_g c_{v,g}\right] dV_i}.
\label{eq:newton}
\end{equation}
The correction is applied a fixed three times with no convergence test, and the
result is clipped to the temperature range spanned by the property tables.

\subsection{Time Step}

The conduction fluxes and the decomposition rates are both evaluated at the
temperature at the start of the step, so the scheme is explicit. An explicit
scheme for a diffusion problem is stable only below a step size that scales with
the square of the mesh spacing~\cite{morton2005pde}. The step size is
set once at startup from that stability estimate,
\begin{equation}
\Delta t = \frac{1}{\mathrm{CFL}} \cdot
\frac{\Delta x_{\max}^2\, \rho_{v,\mathrm{bulk}}\, c_{p,\mathrm{mix}}(600\,\mathrm{K}, \tau_{\mathrm{init}})}
{2\, k_{\mathrm{mix}}(600\,\mathrm{K}, \tau_{\mathrm{init}})} .
\label{eq:dt}
\end{equation}
Three properties of Equation~\eqref{eq:dt} are worth stating plainly. The
parameter $\mathrm{CFL}$ appears in the denominator, so it is a safety divisor
rather than a Courant number, and a larger value gives a smaller step. The
material properties are frozen at $600$~K and at the fully virgin state. The step
is computed once and never adapted, and the only per-step modification is a
clipping so that the step lands exactly on the final time or on the next
breakpoint in a boundary-condition table.

Because $\Delta t \propto \Delta x^2$, the step count for a fixed physical
duration grows quadratically with mesh resolution. This is the cost the surrogate
is intended to avoid.

\section{Pyrolysis}

The solid is represented as a set of $N_s$ species that decompose independently.
Each species $k$ has an intrinsic virgin density $\rho_{v,k}$, an intrinsic char
density $\rho_{c,k}$, a pre-exponential factor $A_k$, a reaction order $\psi_k$,
an activation temperature $E_k$, and a mass-fraction weight $\gamma_k$. This
multi-species Arrhenius description of resin decomposition follows the classical
kinetics measured for phenolic and nylon-phenolic
composites~\cite{goldstein1969kinetics}. The material used here is
TACOT~\cite{lachaud2012tacot}, a reference composite defined for code comparison,
with three species of which two are reactive and one is inert.

The decomposition state of species $k$ is the normalized reactive mass fraction
\begin{equation}
W_k = \mathrm{clip}\!\left(\frac{\rho_{i,k} - \rho_{c,k}}{\rho_{v,k}},\, 0,\, 1\right),
\end{equation}
which obeys the rate law
\begin{equation}
\frac{dW_k}{dt} = -A_k\, W_k^{\psi_k}\, \exp\!\left(-\frac{E_k}{T}\right).
\label{eq:arrhenius}
\end{equation}
Note that $E_k$ is stored already divided by the universal gas constant, so it
carries units of kelvin and no gas constant appears in
Equation~\eqref{eq:arrhenius}.

Equation~\eqref{eq:arrhenius} is not integrated numerically. It is solved in
closed form over the step at frozen temperature,
\begin{equation}
W^{n+1}_k =
\begin{cases}
W_k \exp\!\left[-A_k \Delta t\, e^{-E_k/T}\right], & \psi_k = 1, \\[8pt]
\left[\max\!\left(0,\; W_k^{1-\psi_k} - \Delta t\,(1-\psi_k) A_k e^{-E_k/T}\right)\right]^{\frac{1}{1-\psi_k}},
& \psi_k \neq 1 .
\end{cases}
\label{eq:arrhenius_closed}
\end{equation}
The species density is then recovered as
$\rho^{n+1}_{i,k} = \rho_{c,k} + \rho_{v,k} W^{n+1}_k$, and the bulk density and
the volumetric gas source follow by summation over species,
\begin{equation}
\rho = \sum_{k=1}^{N_s} (1-\varphi_v)\,\gamma_k\, \rho_{i,k},
\qquad
\dot\omega_g = \sum_{k=1}^{N_s} (1-\varphi_v)\,\gamma_k\,
\frac{\max\!\left(0,\, \rho^{n}_{i,k} - \rho^{n+1}_{i,k}\right)}{\Delta t}.
\label{eq:gas_source}
\end{equation}
The mass lost by the solid becomes the pyrolysis gas that the transport equations
of Section~\ref{sec:gas} carry out through the char.

Equation~\eqref{eq:arrhenius_closed} has no activation threshold. A minimum
reaction temperature is present in the material file and is read by the solver,
but it is never applied, so decomposition proceeds at the Arrhenius rate at all
temperatures.

\section{Char State and Material Properties}
\label{sec:tau}

Two distinct measures of decomposition are used, and they are not
interchangeable.

The first is the char-progress variable $\tau$, which drives the thermal
properties. It equals one in virgin material and zero in fully charred material.
It is computed after removing the inert species, whose density never changes and
would otherwise compress the range,
\begin{equation}
\rho_{\mathrm{inert}} = \rho_{v,\mathrm{bulk}}
\frac{\sum_{k \in \mathcal{I}} \rho_{v,k}}{\sum_{k} \rho_{v,k}},
\qquad
\mathcal{I} = \left\{k : \rho_{v,k} = \rho_{c,k}\right\},
\end{equation}
\begin{equation}
\tau = \mathrm{clip}\!\left[
\frac{1 - \rho_{cr}/\rho_r}{1 - \rho_{cr}/\rho_{vr}},\, 0,\, 1\right],
\label{eq:tau}
\end{equation}
where $\rho_{vr}$, $\rho_{cr}$, and $\rho_r$ are the virgin, char, and current
bulk densities with the inert contribution subtracted.

The second is the linear char fraction $\beta$, which drives porosity and
permeability. It is the degree-of-char variable used in the material-response
literature~\cite{amar2016char},
\begin{equation}
\beta = \mathrm{clip}\!\left[
\frac{\rho_{v,\mathrm{bulk}} - \rho}{\rho_{v,\mathrm{bulk}} - \rho_{c,\mathrm{bulk}}},\, 0,\, 1\right].
\label{eq:beta}
\end{equation}
Porosity is interpolated linearly in $\beta$ between the virgin and char values,
\begin{equation}
\varphi = \mathrm{clip}\left[\varphi_v + (\varphi_c - \varphi_v)\beta,\;
\varphi_{\min},\, \varphi_{\max}\right].
\end{equation}

Thermal conductivity and specific heat are blended linearly in $\tau$ between
separately tabulated virgin and char curves, each of which is a function of
temperature. Tabulating virgin and char properties separately and interpolating
on char state is the standard treatment for a low-density carbon-phenolic
ablator~\cite{milos2010pica},
\begin{equation}
k_{\mathrm{mix}}(T,\tau) = \tau\, k_v(T) + (1-\tau)\, k_c(T),
\qquad
c_{p,\mathrm{mix}}(T,\tau) = \tau\, c_{p,v}(T) + (1-\tau)\, c_{p,c}(T).
\label{eq:mixing}
\end{equation}

The sensible enthalpy is built from the same tables but is referenced so that it
vanishes at the reference temperature for every char state,
\begin{equation}
h_{s,\mathrm{form}}(\tau) = \tau\, e_v(T_{\mathrm{ref}}) + (1-\tau)\, e_c(T_{\mathrm{ref}}),
\end{equation}
\begin{equation}
h_{s}(T,\tau) = \tau\, e_v(T) + (1-\tau)\, e_c(T) - h_{s,\mathrm{form}}(\tau),
\label{eq:hs}
\end{equation}
where $e_v$ and $e_c$ are the integrated heat capacities
$e(T) = \int^{T} c_p(T')\,dT'$ of the virgin and char material. Equation
\eqref{eq:hs} gives $h_s(T_{\mathrm{ref}}, \tau) = 0$ for all $\tau$, which is
what makes the formation term of Equation~\eqref{eq:echem} the sole carrier of
the pyrolysis energy.

The surface emissivity is blended in $\tau$ in the same way,
$\varepsilon = \tau \varepsilon_v + (1-\tau)\varepsilon_c$.

\section{Pyrolysis-Gas Transport}
\label{sec:gas}

The gas generated by Equation~\eqref{eq:gas_source} percolates out through the
porous char toward the heated surface. The solver treats this as quasi-steady
Darcy flow, which is the usual momentum closure for pyrolysis gas in a charring
ablator~\cite{weng2014pyrolysisgas}. There is no gas momentum equation and no
transient gas-mass matrix, so the pressure field is re-solved from scratch at
every step.

Darcy's law, which follows from volume-averaging creeping flow through a porous
medium~\cite{whitaker1986darcy}, relates the gas mass flux to the pressure
gradient,
\begin{equation}
\dot m'' = -\frac{K_0 \rho_g}{\mu} \frac{\partial p_g}{\partial x},
\label{eq:darcy}
\end{equation}
which discretizes to a conductance form on the mesh,
\begin{equation}
\dot m_i = -C_i \left(p_{g,i} - p_{g,i-1}\right),
\qquad
C_i = \frac{K_{0,i}\, \rho_g(T_{f,i}, p_{f,i})\, dA_i}
{\mu(T_{f,i}) \cdot \tfrac12\left(\Delta x_{i-1} + \Delta x_i\right)} .
\label{eq:darcy_discrete}
\end{equation}
The permeability is blended between its virgin and char values using the linear
char fraction of Equation~\eqref{eq:beta} rather than $\tau$,
\begin{equation}
K_0 = (1 - \beta_f) K_{0,v} + \beta_f K_{0,c} .
\end{equation}
No Klinkenberg slip correction or Forchheimer inertial term is included. Departures
from Darcy behavior can matter for pyrolysis gas at low pressure and high
flux~\cite{martin2010nondarcian}, and neglecting them is a modeling choice
inherited from the reference implementation.

Gas mass conservation closes the system,
\begin{equation}
\frac{\partial (\varphi \rho_g)}{\partial t} + \frac{\partial \dot m''}{\partial x}
= \dot\omega_g ,
\end{equation}
which is discretized with the storage term moved to the right-hand side so that
the remaining balance is steady,
\begin{equation}
\dot m_{i+1} - \dot m_i = \dot\omega_{\mathrm{net},i}\, dV_i,
\qquad
\dot\omega_{\mathrm{net},i} = \dot\omega_{g,i}
- \frac{\left(\varphi \rho_g\right)^{n}_i - \left(\varphi \rho_g\right)^{n-1}_i}{\Delta t}.
\label{eq:gas_mass}
\end{equation}

For the boundary configuration used in all cases considered here, the back face
is sealed and the heated surface is open to the ambient pressure. Equation
\eqref{eq:gas_mass} then admits a direct march with no linear solve,
\begin{equation}
\dot m_1 = 0, \qquad
\dot m_{i+1} = \dot m_i + \dot\omega_{\mathrm{net},i}\, dV_i,
\end{equation}
\begin{equation}
p_{g,N+1} = p_{\mathrm{amb}}(t), \qquad
p_{g,i-1} = p_{g,i} + \frac{\dot m_i}{C_i} .
\end{equation}
For other boundary combinations the same conductances assemble a tridiagonal
system in $p_g$ that is solved directly. The ambient pressure is time-varying and
is taken from the aerothermal table.

The gas is treated as ideal with tabulated, temperature-dependent properties,
\begin{equation}
\rho_g(T, p) = \frac{p}{R(T)\,\max(T, 250\,\mathrm{K})},
\qquad
c_{v,g} = c_{p,g}(T) - R(T),
\end{equation}
where the floor on temperature in the denominator is a numerical guard carried
over from the reference implementation.

\section{Boundary Conditions}

Boundary conditions are imposed by setting the ghost-cell temperatures so that
the reconstructed face gradient carries the prescribed flux.

\subsection{Aerothermal Surface}

At the heated surface the net flux into the wall is the sum of convective
heating, incident radiation, and the enthalpy carried by pyrolysis gas leaving the
surface, less the energy re-radiated away. This is the surface energy balance of
the charring-material ablation analysis~\cite{moyer1968cma}, and it is what
couples the in-depth solution to the external
aerothermodynamics~\cite{martin2015coupled},
\begin{equation}
q_w = \underbrace{\left(\rho_e u_e C_H\right) \Phi \left(h_r - h_w\right)}_{\text{convective}}
    + \underbrace{q_{\mathrm{rad}}}_{\text{incident}}
    + \underbrace{\frac{\dot m_{N+1}}{dA_{N+1}}\left(h_g^{\mathrm{abs}}(T_w) - h_w\right)}_{\text{pyrolysis gas}}
    - \underbrace{\varepsilon \sigma \left(T_w^4 - T_{\mathrm{ext}}^4\right)}_{\text{re-radiation}} ,
\label{eq:surface_balance}
\end{equation}
where $\rho_e u_e C_H$ is the heat-transfer coefficient, $h_r$ the recovery
enthalpy, $h_w$ the wall enthalpy, $\sigma$ the Stefan-Boltzmann constant, and
$T_w$ the wall temperature. The coefficients $\rho_e u_e C_H$, $h_r$,
$q_{\mathrm{rad}}$, and $T_{\mathrm{ext}}$ are interpolated in time from a
boundary-layer table.

The factor $\Phi$ is the blowing correction, which accounts for the way gas
leaving the surface thickens the boundary layer and reduces the convective heat
transfer~\cite{moyer1968cma}. The reference implementation evaluates it in the
exponential form
\begin{equation}
\varphi_b = \frac{2\lambda\, \dot m''_w}{\left(\rho_e u_e C_H\right)_0},
\qquad
\Phi = \frac{\varphi_b}{e^{\varphi_b} - 1} ,
\label{eq:blowing}
\end{equation}
with $\lambda$ a tabulated blowing parameter. The limit $\Phi \to 1$ as
$\varphi_b \to 0$ recovers the unblown heat transfer, and the implementation uses
a series expansion in that limit to avoid the removable singularity.

The wall enthalpy $h_w$ is not computed from a chemistry model at run time. It is
read from precomputed surface-thermochemistry tables indexed by the wall
temperature and by the dimensionless blowing rate. Such tables are generated by an
equilibrium surface-thermochemistry solution~\cite{kendall1968ace}, and the
underlying mass and energy balance that defines $B'$ has since been derived from
first principles~\cite{padovan2024bprime},
\begin{equation}
B'_g = \frac{\dot m_{N+1} / dA_{N+1}}{\left(\rho_e u_e C_H\right)_0 \Phi},
\qquad
h_w = h_w\!\left(T_w,\, B'_g\right),
\label{eq:bprime}
\end{equation}
interpolated bilinearly and clamped outside the tabulated range. Only the
tables at one ambient pressure are read, so there is no pressure interpolation of
the surface chemistry.

Finally the ghost temperature is set so the face gradient carries $q_w$,
\begin{equation}
T_{N+1} = T_N + \frac{q_w}{k_{\mathrm{mix}}}\, \Delta x_N .
\end{equation}

For the gas-off case the surface condition is instead a prescribed wall-temperature
history imposed as a Dirichlet condition, $T_{N+1} = 2 T_w(t) - T_N$.

\subsection{Back Face}

The back face is adiabatic in all cases considered here. The applied flux is zero
and the back-face re-radiation term is disabled in the reference implementation,
so the ghost temperature simply mirrors the interior,
\begin{equation}
T_0 = T_1 - \frac{q_{w,0}}{k_{\mathrm{mix}}}\, \Delta x_1, \qquad q_{w,0} = 0 .
\end{equation}

\section{Step Ordering}

The order of operations within one step is fixed and is worth recording, because
the surrogate reproduces the same sequence. The step size is selected, the ambient
pressure is evaluated, the boundary conditions are applied to set the ghost cells,
the species are advanced by Equation~\eqref{eq:arrhenius_closed}, the gas system of
Equation~\eqref{eq:gas_mass} is solved, the face conductivities of
Equation~\eqref{eq:kface} are formed, the cell energies are updated by
Equation~\eqref{eq:energy_update}, and the temperature is recovered by
Equation~\eqref{eq:newton}. The surface boundary condition at step $n$ uses the
gas mass flux computed at step $n-1$, which is an explicit lag inherited from the
reference implementation.

\section{Scope and Limitations}

Several features that appear in the broader charring-ablator literature are
absent from the solver used here, and the surrogate therefore inherits their
absence. They are listed in Table~\ref{tab:limitations}.

\begin{table}[H]
\centering
\caption{Physics present in the reference implementation but not exercised in
this work, or absent entirely.}
\label{tab:limitations}
\begin{tabular}{ll}
\toprule
Feature & Status \\
\midrule
In-depth radiative transport & Not implemented. Radiation enters only at the surface. \\
Surface recession & Not implemented. The mesh is fixed. \\
Species activation threshold & Read from the material file, never applied. \\
Pressure-dependent surface chemistry & Only one ambient-pressure table is read. \\
Gas momentum and inertia & Quasi-steady Darcy only. \\
Adaptive time stepping & Step size fixed at startup. \\
Non-uniform mesh & Uniform spacing only. \\
\bottomrule
\end{tabular}
\end{table}

The one-dimensional restriction is the most consequential of these. Heat and
pyrolysis gas are resolved only through the thickness, so lateral conduction and
in-plane gas motion are absent. Two-dimensional and fully three-dimensional
material-response codes exist~\cite{chen2001titan,amar2016char,lachaud2014pato},
and extending the surrogate to their meshes is left to future work. The choice is
deliberate here. A one-dimensional solver is cheap enough to generate the large
training sets this work requires, which makes it a tractable testbed for the
surrogate method.

\bibliographystyle{unsrt}
\bibliography{heat_gnn_refs}

\end{document}
